% Verbatim copy of /Users/alingalatan/workspaces/grothe/handoff/manuscript/legacy/ch25-parts/60-anisotropic-cap.tex (file dated 2026-10-25). Only this line was added.
\subsection{A sharp two-disk majorant and a larger original-response region}
\label{sec:anisotropic-response-cap}
\paragraph{Provenance (IDE note, 21 October).} Lower-bound GPT two-disk package (\code{conjK-anisotropic-response-20261021.zip}, archived in
\pathref{gpt-reviews/anisotropic-response-20261021/}). IDE replays \code{python3 -I verify\_all.py} and \code{--fine} (standard library, unrelated directory): ALL CLAIMS
PASS (1{,}897 and 5{,}651 panels; 5.3\,s and 15.5\,s), including an exact sum-of-squares identity for the two-disk majorant. IDE check: random search plus local
maximisation of the pointwise two-disk inequality over complex $h,k$ with $|h\pm k|\le1$, complex $L$ and $\lambda>0$ gives maximum $0$ (sharp, never violated).
Bookkeeping: the inverse is real-linear (the vector $(r_A,\overline{r_B})$ and its conjugate term are kept); the $\|Pk\|^2$ square is used by the rank-one corrections
and must not be subtracted again; $(a,b)=(0.4541,0.15245)$ is a dual choice, not a profile change. The older operator theorem stays valid on its slightly wider
$t$-interval; take minima of valid bounds. Not a raw-K region, not a Lean formalisation.


\paragraph{Scope and status.}
[P] means analytic proof and [C] exact directed computation.
The statements concern the original bounded-response problem, in every
finite dimension with all visible sectors and arbitrary shared drivers.
No fourth- or sixth-moment input, raw Conjecture K, stationarity, or fresh
Gaussian buffer is assumed. They enlarge the preceding small-remote-variance
region; they do not cover the complete complement or give a new global
Grothendieck endpoint.

Fix $\epsilon=1891/2500$, $p=1/\epsilon$, and write
\[
 h=(f_++f_-)/2,\quad k=(f_+-f_-)/2,\quad |f_\pm|\le1,
 \quad Ph=tZ,\quad F=\epsilon Qk,
\]
with Gaussian Wick projections $P=P_{1,0}$ and $Q=P_{2,1}$. The visible
standard proper complex Gaussian $Z$ is independent of the remote Gaussian
space $X$. Set $T=|Z|^2$ and decompose
\[
 F=q+(T-1)r_A+Z^2r_B/\sqrt2+ZY+\bar ZW+\nu H(Z),
 \qquad H(Z)=Z(T-2)/\sqrt2.
\]
The remote bidegrees are $(2,1),(1,0),(0,1),(1,1),(2,0)$ in that order.
The remote fields need not be independent of one another. Let
\[
 e=\|q\|_2^2,\quad A_2=\|r_A\|_2^2,\quad B_2=\|r_B\|_2^2,
 \quad Y_R=\operatorname{Re}Y,\quad Y_I=\operatorname{Im}Y.
\]
The exact original objective is
\[
 \mathcal N=t^2+p\|F\|_2^2-\|Pk\|_2^2-\epsilon\|Qh\|_2^2.
\]

\subsubsection{The new regional theorem [P,C]}
For every actual response,
\begin{equation}\label{eq:aniso-rational}
\boxed{\begin{aligned}
 \mathcal N\le{}&t^2-\frac{759}{500}t+\frac{1272881}{10^6}
                   +\frac{27}{50}e-\frac15A_2-\frac9{10}B_2\\
 &-\frac{1219}{1000}\|Y_R\|^2-\frac1{4000}\|Y_I\|^2
       -\frac{417}{1000}\|W\|^2\\
 &-\frac{89}{40}(\operatorname{Re}\nu)^2
       -\frac{163}{250}(\operatorname{Im}\nu)^2.
\end{aligned}}
\end{equation}
Consequently
\begin{equation}\label{eq:aniso-region}
 \boxed{\begin{gathered}
 \frac{667}{1000}\le t\le\frac{851}{1000},\qquad
 e\le\frac{37}{2500}=0.0148\\
 \Longrightarrow\quad
 \mathcal N\le\frac{89157}{125000}=0.713256<\frac{500}{701}.
 \end{gathered}}
\end{equation}
In particular the entire central interval $[0.84,0.851]$ is included.
The variance cutoff increases from $0.0113$ to $0.0148$, by more than
$30\%$, with every sector still unrestricted. This is a regional statement
about the unnormalized variance of the actual projection
$q=\mathbb E[F\mid X]$.

The A/B penalty in~\eqref{eq:aniso-rational} can be strengthened to
\begin{equation}\label{eq:aniso-AB}
 \mathcal P_{AB}=\frac{57}{125}A_2+
 \frac{43}{50}\operatorname{Re}\langle r_A,\overline{r_B}\rangle
       +\frac{1623}{1000}B_2.
\end{equation}
The corresponding real symmetric matrix dominates both
$\operatorname{diag}(1/5,9/10)$ and $(31/100)I$. Thus this is a positive
penalty at every relative phase; its off-diagonal term must not be discarded
with an arbitrary sign.

Define
\[
\begin{aligned}
 \mathcal E_{\rm adj}={}&e-\frac{10}{27}A_2-\frac53B_2
 -\frac{1219}{540}\|Y_R\|^2-\frac1{2160}\|Y_I\|^2
 -\frac{139}{180}\|W\|^2\\
 &-\frac{445}{108}(\operatorname{Re}\nu)^2
 -\frac{163}{135}(\operatorname{Im}\nu)^2.
\end{aligned}
\]
On the central interval, the larger condition
\begin{equation}\label{eq:aniso-adjusted}
 \mathcal E_{\rm adj}\le\frac{37}{2500}
       +\frac3{10}\left(\frac{851}{1000}-t\right)
\end{equation}
also gives the cap in~\eqref{eq:aniso-region}. One may replace the two
A/B terms in $\mathcal E_{\rm adj}$ by $-(50/27)\mathcal P_{AB}$.

\subsubsection{The sharp pointwise inequality [P]}
The previous energy relaxation used only $|h|^2+|k|^2\le1$.
The two disks contain more information. For any $L\in\mathbb C$,
$\lambda>0$, and $|h\pm k|\le1$,
\begin{equation}\label{eq:aniso-disk}
 \boxed{\quad
 2\operatorname{Re}(\overline Lh)+\lambda|k|^2
 +\frac{\{\operatorname{Re}(\overline Lk)\}^2}{\lambda}
 \le\lambda+\frac{|L|^2}{\lambda}.
 \quad}
\end{equation}
For $L\ne0$, this adds the term
$(|L|^2/\lambda)\{\operatorname{Re}(\overline Lk/|L|)\}^2$.
Its coefficient is sharp with the other terms fixed: $h=0,k=L/|L|$
attains equality.

Here is an elementary proof. First take $f_\pm=h\pm k$ on the unit circle,
rotate $L$ to $l=|L|\ge0$, and write
$h=e^{i\phi}\cos\psi$, $k=ie^{i\phi}\sin\psi$.
Put $a_0=\cos\phi$, $c_0=\cos\psi$ and $x=l/\lambda$.
The difference between right and left sides of~\eqref{eq:aniso-disk},
divided by $\lambda$, equals
\[
\begin{aligned}
 &1+x^2-2xa_0c_0-(1-c_0^2)-x^2(1-a_0^2)(1-c_0^2)\\
 &\hspace{15mm}=(c_0-xa_0)^2+x^2c_0^2(1-a_0^2)\ge0.
\end{aligned}
\]
For disks, write each $f_\pm$ as the mean of two antipodal unit-circle
values. The left side is a convex quadratic function of the pair
$(f_+,f_-)$, so Jensen extends the inequality. This does not assume
unimodularity of the original responses. The displayed polynomial identity
is also checked exactly by the replay.

\subsubsection{The weighted real Hilbert-space reduction [P]}
Choose $a,b>0$, put
\[
 L_0=Z(a+bT),\quad \ell_0=a+2b,\quad
 \mathcal B=bZ(T-2),\quad l(T)=\sqrt T(a+bT),
\]
and choose a positive radial $\lambda(T)$. Applying~\eqref{eq:aniso-disk}
pointwise leads to the quadratic form
\[
 \mathfrak m(k)=\mathbb E\left[\lambda|k|^2+
 \frac{l^2}{\lambda}\{\operatorname{Re}(e^{-i\arg Z}k)\}^2\right].
\]
At $Z=0$ choose the phase arbitrarily. The inverse of its pointwise real
linear operator is
\[
 M^{-1}u=c(T)u+d(T)e^{2i\arg Z}\bar u,
\]
where
\[
 c(T)=\frac12\left(\lambda^{-1}+\frac{\lambda}{\lambda^2+l^2}\right),
 \qquad
 d(T)=\frac12\left(\frac{\lambda}{\lambda^2+l^2}-\lambda^{-1}\right).
\]
In particular one must keep the conjugate term; treating the operator as
complex linear would give incorrect A/B cross terms.

Define, for $0\le j\le3$,
\[
 J_j=\mathbb E\frac{T^j}{\lambda},\quad
 V_j=\mathbb E\frac{T^j\lambda}{\lambda^2+l^2},\quad
 C_j=(J_j+V_j)/2,\quad D_j=(V_j-J_j)/2,
\]
and for any four-vector $U$ put
\[
 K_H(U)=\frac{U_3-4U_2+4U_1}{2}
             -\frac{(U_2-2U_1)^2}{2(1+U_1)}.
\]
The coupled A/B block is the real matrix
\begin{equation}\label{eq:aniso-gram}
 K=\begin{pmatrix}
 C_2-2C_1+C_0&(D_2-D_1)/\sqrt2\\
 (D_2-D_1)/\sqrt2&C_2/2
 \end{pmatrix}
 -\frac1{1+C_0}
 \begin{pmatrix}C_1-C_0\\D_1/\sqrt2\end{pmatrix}
 \begin{pmatrix}C_1-C_0&D_1/\sqrt2\end{pmatrix}.
\end{equation}
If all inverse quantities exist and $K$ is positive definite, then
\begin{equation}\label{eq:aniso-coercivity}
\boxed{\begin{aligned}
 \mathfrak m(k)+\|Pk\|^2\ge p^2\bigg\{&e/C_0+
 \left\langle\binom{r_A}{\overline r_B},
 K^{-1}\binom{r_A}{\overline r_B}\right\rangle\\
 &+\|Y_R\|^2/V_1+\|Y_I\|^2/J_1+\|W\|^2/C_1\\
 &+(\operatorname{Re}\nu)^2/K_H(V)
   +(\operatorname{Im}\nu)^2/K_H(J)\bigg\}.
\end{aligned}}
\end{equation}
The matrix quadratic form uses the real part of the complex inner product.

For clarity, the complete block bookkeeping is as follows. Decompose the
remote Gaussian space into Hermite chaoses and group each chaos with its
conjugate, since $M^{-1}$ conjugates. The selected moments in different
listed groups are orthogonal for the inverse weighted form.
\[
\begin{array}{c|c|c}
 \text{selected cubic component}&\text{inverse Gram value}
       &\text{retained first-chaos correction}\\\hline
 q&C_0&\text{none}\\
 (r_A,\overline r_B)&K&\text{remote }P_{1,0}k\\
 Y_R,\ Y_I&V_1,\ J_1&\text{none}\\
 W&C_1&\text{none}\\
 \operatorname{Re}\nu,\ \operatorname{Im}\nu&K_H(V),\ K_H(J)
       &\text{visible real/imaginary }Pk
\end{array}
\]
For one remote linear coordinate, the selected visible functions are
$X(T-1)$, $\bar XZ^2/\sqrt2$, and the penalized first-chaos function $X$.
The raw A/B Gram matrix is the first matrix in~\eqref{eq:aniso-gram};
its pairings with $X$ are $(C_1-C_0,D_1/\sqrt2)$ and the inverse norm
of $X$ is $C_0$. Adding $|\langle X,k\rangle|^2$ subtracts the stated
rank-one term. The real and imaginary coordinate blocks agree after
conjugating the B coefficient.

The remote $(1,1)$ chaos has a real orthonormal basis of real-valued
Hermitian quadratic polynomials. On $Z$ times that basis the real inverse
weight is $\lambda/(\lambda^2+l^2)$ and the imaginary one is
$\lambda^{-1}$, giving $V_1,J_1$. The remote $(2,0)$ block gives $C_1$;
its conjugate term has the wrong visible frequency to pair with a selected
cubic. The visible diagonal gives the two rank-one corrections $K_H$.
The remote cubic has no selected conjugate-frequency partner, giving
$C_0$. Weighted Bessel, or minimization of the weighted norm subject to
these moments, proves~\eqref{eq:aniso-coercivity}. Every unlisted component
is allowed in that minimization. This proves the bound in arbitrary
dimension, without independence between any remote fields.

More explicitly, let $A=M+P$ on the real Hilbert space with inner product
$\operatorname{Re}\mathbb E(\bar u v)$, and let $u_i$ be a real basis of
the selected cubic moments. If
$G_{ij}=(u_i,A^{-1}u_j)_{\mathbb R}$ and
$y_i=(u_i,k)_{\mathbb R}$, then, for every real vector $z$,
\[
 (k,Ak)_{\mathbb R}\ge 2z^{\mathsf T}y-z^{\mathsf T}Gz.
\]
Maximizing over $z$ gives $y^{\mathsf T}G^{-1}y$. Here $y$ consists of
$p$ times the coefficients of $F=\epsilon Qk$. The finite-rank identity
\[
 (M+P)^{-1}=M^{-1}-M^{-1}P(I+PM^{-1}P)^{-1}PM^{-1}
\]
on the range of $P$ gives precisely the corrections above. For the
certified weight, $\lambda\ge\sqrt{s_0}>0$ and all selected polynomial
moments are finite. Thus the argument can equivalently be read as a
finite-dimensional dual minimization, without any boundedness assumption
on multiplication by $\lambda$.

\subsubsection{Completion of the original objective [P]}
The identity
\[
 -\epsilon\|Qh\|^2
 =2\operatorname{Re}\langle\mathcal B,h\rangle+2pb^2
                  -\epsilon\|Qh+p\mathcal B\|^2
\]
and $Ph=tZ$ convert the original objective into a support problem for
$L_0$. Use~\eqref{eq:aniso-disk} and then~\eqref{eq:aniso-coercivity}.
With
\[
 C_\lambda=2pb^2+\mathbb E(\lambda+l^2/\lambda),\qquad
 \Delta_{AB}=p^2K^{-1}-pI,
\]
the result is
\begin{equation}\label{eq:aniso-general}
\begin{aligned}
 \mathcal N\le{}&t^2-2\ell_0t+C_\lambda+(p-p^2/C_0)e
 -\left\langle\binom{r_A}{\overline r_B},\Delta_{AB}
                         \binom{r_A}{\overline r_B}\right\rangle\\
 &-(p^2/V_1-p)\|Y_R\|^2-(p^2/J_1-p)\|Y_I\|^2
       -(p^2/C_1-p)\|W\|^2\\
 &-(p^2/K_H(V)-p)(\operatorname{Re}\nu)^2
       -(p^2/K_H(J)-p)(\operatorname{Im}\nu)^2\\
 &-\epsilon\|Qh+p\mathcal B\|^2.
\end{aligned}
\end{equation}
All penalties belong to one derivation. The $Qh$ square is still available;
it was not used in the weighted operator calculation. The parameters
$a,b$ are free parameters of this original-response dual, not a change
to the peak profile or the witness in Conjecture K.

\subsubsection{Exact numerical certificate [C]}
Take
\[
 a=\frac{4541}{10000},\quad b=\frac{3049}{20000},\quad
 s_0=\frac{167}{50000},\quad d=\frac{11}{500},\qquad
 \lambda(T)=\sqrt{T(a+bT)^2+s_0+dT}.
\]
Then $2\ell_0=759/500$. The replay proves
\[
 C_\lambda<1.272881,\quad 0<p-p^2/C_0<0.54,\quad
 \Delta_{AB}\succeq
 \begin{pmatrix}0.456&0.430\\0.430&1.623\end{pmatrix},
\]
and the five remaining penalties in order $(Y_R,Y_I,W,\nu_R,\nu_I)$
are respectively greater than
\[
 1.219,\qquad0.00025,\qquad0.417,\qquad2.225,\qquad0.652.
\]
All decimals in these certified comparisons are exact rationals.
For the A/B simplification, the matrix minus
$\operatorname{diag}(0.2,0.9)$ has positive diagonal and determinant
\[
 (0.456-0.2)(1.623-0.9)-0.430^2=0.000188>0.
\]
Subtracting $0.31I$ also leaves a positive definite matrix.
This proves~\eqref{eq:aniso-rational} and~\eqref{eq:aniso-AB}.
The two endpoints $0.667$ and $0.851$ sum to $2\ell_0$; convexity
in $t$ proves~\eqref{eq:aniso-region}. The margin is exactly
$943/87625000>10^{-5}$. Substituting~\eqref{eq:aniso-adjusted}
gives a quadratic nondecreasing on the central interval, with the same
upper-endpoint value.

Only eight one-dimensional integrals are needed. Write
\[
 P(T)=T(a+bT)^2+s_0+dT,\qquad
 Q_0(T)=2T(a+bT)^2+s_0+dT.
\]
The two integrands, apart from $T^je^{-T}$, are $P^{-1/2}$ and
$\sqrt P/Q_0$. The expectation in $C_\lambda$ is just
\[
 s_0J_0+(2a^2+d)J_1+4abJ_2+2b^2J_3.
\]
The verifier uses fixed-point integer intervals at scale $10^{40}$ and
adaptive Simpson panels on $[0,40]$. Fourth derivatives are enclosed over
the entire panel by Taylor-coefficient recurrences from
$2Pu'+P'u=0$ and $Q_0v=1$. The panel error is at most
$h^5\sup|f^{(4)}|/2880$, or $h^5/120$ times the enclosure of the
normalized fourth coefficient $f^{(4)}/4!$. Exponentials have explicit positive
Taylor remainders, and square roots use integer square roots.
Both tails obey
\[
 0\le\int_{40}^\infty T^je^{-T}P^{-1/2}\,dT,
 \ \int_{40}^\infty T^je^{-T}\frac{\sqrt P}{Q_0}\,dT
 \le\frac{j!e^{-40}}b\sum_{k=0}^j\frac{40^k}{k!},
\]
since $Q_0\ge P\ge b^2T^3$. No stored integral or floating quadrature
is accepted as proof input. The fine replay decreases the integration
error target by a factor of $100$.

\subsubsection{The remaining complement}
Any central-row response with $\mathcal N\ge500/701$ must now satisfy
\[
 \mathcal E_{\rm adj}\ge
 \frac{500/701-t^2+(759/500)t-1272881/10^6}{27/50}.
\]
The right side is approximately $0.0183440$ at $t=0.84$,
$0.0166059$ at $t=0.8456$, and $0.0148199$ at $t=0.851$.
The full orientation penalty~\eqref{eq:aniso-AB} can strengthen this
condition. The candidate must also evade every existing cumulant/support
test and the earlier operator certificates. These necessary conditions
have not been shown jointly impossible. The intermediate region and the
full same-$\epsilon$ row cover remain open; the global certified interval
is unchanged.

